% ============================================================================
% Fichier  : front/02_mode_emploi.tex
% Titre    : Mode d'emploi du cours
% Auteur   : MINKA MI NGUIDJOÏ Thierry Emmanuel
% Version  : 0.1 -- novembre 2026
% ============================================================================
% 8 jours de 4 h ; les 2 dernières heures du J08 sont offertes (30 h
% effectives). Fiches et labs hors de ces heures.
% ============================================================================

\chapter*{Mode d'emploi du cours}
\addcontentsline{toc}{chapter}{Mode d'emploi du cours}
\markboth{Mode d'emploi}{Mode d'emploi}

\section*{Structure générale}

Le cours compte un prologue et six parties, soit quinze chapitres. Il
est donné en huit journées de quatre heures~; les deux dernières heures de
la huitième journée sont offertes aux auditeurs, ce qui fait trente heures
effectives. Les fiches d'auto-évaluation et les laboratoires de
60~minutes de NDAP SUKULU sont des travaux à faire hors de ces heures.

\begin{tableaucours}{L{1.6cm}L{4.4cm}L{4.4cm}L{1.8cm}}{%
    \tableauHeader{Partie} &
    \tableauHeader{Thème} &
    \tableauHeader{Compétences visées} &
    \tableauHeader{Niveau}
}
Prologue & Le paiement du 28 décembre & Voir ce que le dossier papier ne
montre pas & \niveaubase \\
I & Fondements et contrôle interne & Situer l'audit d'un SI public,
décrire et évaluer un système de contrôle interne & \niveaubase{} à
\niveaumoyen \\
II & Cadre institutionnel et normatif & Savoir qui contrôle quoi, choisir
ses normes et ses critères & \niveaubase \\
III & Démarche de mission & Préparer et décomposer, coter \OUI{}~/
\NON{}~/ \NEK{}, mesurer $\mu$, $U$, $H$, conclure, rapporter, suivre les
cycles & \niveaumoyen \\
IV & Domaines d'audit & Évaluer les accès, l'exploitation, les
applications, les données et les projets & \niveaumoyen{} à
\niveauexpert \\
V & Synthèse & Conduire une mission complète en équipe & \niveauexpert \\
VI & Ouverture sur la recherche & Situer l'auditologie, formuler un sujet
de mémoire falsifiable & \niveauexpert \\
\end{tableaucours}

\section*{Les huit journées}

\begin{tableaucours}{L{1.4cm}L{8.6cm}L{2.2cm}}{%
    \tableauHeader{Jour} &
    \tableauHeader{Contenu} &
    \tableauHeader{Partie}
}
J01 & Prologue~; contrôler et auditer un SI public
(Ch.~\ref{chap:audit-public}) & Prologue, I \\
J02 & Le système de contrôle interne (Ch.~\ref{chap:sci})~; évaluer le
système de contrôle interne (Ch.~\ref{chap:eval-sci}) & I \\
J03 & Qui contrôle les SI publics (Ch.~\ref{chap:institutions})~; normes
et référentiels (Ch.~\ref{chap:normes}) & II \\
J04 & La mission~: principes, préparation, décomposition
(Ch.~\ref{chap:preparer})~; collecter la preuve~: \OUI{}, \NON{}, \NEK{}
(Ch.~\ref{chap:preuve}) & III \\
J05 & Mesurer, conclure, rapporter, suivre (Ch.~\ref{chap:constater})~;
accès et habilitations (Ch.~\ref{chap:acces}) & III, IV \\
J06 & Changements, exploitation et continuité
(Ch.~\ref{chap:exploitation})~; contrôles d'application
(Ch.~\ref{chap:applications}) & IV \\
J07 & Analyse de données et fraude (Ch.~\ref{chap:donnees-fraude})~;
gouvernance, projets et prestataires (Ch.~\ref{chap:gouvernance}) & IV \\
J08 & Cas intégrateur en salle (Ch.~\ref{chap:cas-integrateur}) et
ouverture sur la recherche (Ch.~\ref{chap:recherche}), deux heures~; les
deux heures suivantes sont offertes & V, VI \\
\end{tableaucours}

\section*{Les encadrés}

\begin{boxdef}
\textbf{Définition}. Introduit un terme avec sa définition~; les termes
définis sont indexés et repris dans le glossaire (Annexe~\ref{ann:glossaire}).
\end{boxdef}

\begin{boxscenario}{Titre du scénario}
\textbf{Scénario}. Un épisode d'un des trois dossiers du cours~: MINÉQUIR,
CAPNO, e-RECETTES.
\end{boxscenario}

\begin{boxobjet}{Code · Intitulé}
\textbf{Point de contrôle}. Présente un contrôle sous une forme
constante~: objectif, risque couvert, test de conception, test de
fonctionnement, preuve attendue.
\end{boxobjet}

\begin{boxprocedure}
\textbf{Papier de travail}. Un modèle de feuille de mission, prêt à
remplir~: ce que l'auditeur fait et ce que le dossier doit contenir.
\end{boxprocedure}

\begin{boxlecture}{Titre de la lecture}
\textbf{Lecture auditologique}. Des réponses \OUI{}~/ \NON{}~/ \NEK{},
avec la constatation, l'analyse et son niveau de confiance, et, selon le
cas, le décompte des contrôles cotés ou les trois mesures du domaine.

\medskip
\hfill\conformite{3}{8}{0}\quad\lecture{0{,}55}{0{,}30}{1{,}41}
\end{boxlecture}

\begin{boxmetier}{Vérificateur}
\textbf{Application métier}. Traduit le chapitre en gestes concrets pour
l'un des métiers visés~: vérificateur du Contrôle Supérieur de l'État,
juge des comptes, inspecteur, contrôleur financier, auditeur interne.
\end{boxmetier}

\begin{boxoutils}{Titre de l'outil}
\textbf{Outils}. Une formule de tableur, une requête ou une commande, à
refaire sur les données du cas.
\end{boxoutils}

\begin{boxjuris}
Cite un texte de loi ou de règlement, lu dans sa version disponible, puis
en tire la portée pour l'audit.
\end{boxjuris}

\begin{boxalert}
\textbf{Violation typique} et \textbf{règle}~: une erreur fréquente, et ce
qu'il faut faire à la place.
\end{boxalert}

\begin{boxresume}
\textbf{À retenir}. Les points essentiels du chapitre.
\end{boxresume}

\begin{boxexo}
\textbf{Exercice}. Gradués en trois niveaux~: \niveaubase\ fondamental,
\niveaumoyen\ intermédiaire, \niveauexpert\ expert. Les corrigés figurent
dans le livret de correction.
\end{boxexo}

\section*{Conventions}

\begin{itemize}
    \item Chaque contrôle testé reçoit une réponse~: \OUI{} (le contrôle
          existe et fonctionne), \NON{} (il manque ou ne fonctionne pas),
          \NEK{} (connaissance insuffisante pour conclure). Un contrôle sans
          objet est écarté, raison écrite.
    \item L'absence d'une trace que le contrôle devait produire vaut \NON{},
          et non \NEK{}.
    \item $\mu$ est le taux de conformité, $U$ le taux d'incertitude, $H$
          l'entropie en bits (au plus $\log_2 3 \approx 1{,}585$), selon
          l'article A1~\cite{minka2026nek}~; la règle de décision examine $U$
          d'abord, $\mu$ ensuite, $H$ pour contrôler. Les seuils du quadrant
          ($\mu = 0{,}5$~; $U = 0{,}20$) sont des conventions du cours.
    \item Les axiomes de l'auditologie portent la numérotation de l'article
          A0 (Chapitre~\ref{chap:recherche}).
    \item Chaque analyse indique un niveau de confiance~: faible, moyenne ou
          élevée.
    \item Les \termecle{termes clés} apparaissent en gras~; les normes
          s'écrivent \norme{ISSAI~100}, les textes camerounais
          \loicam{décret n°~2013/287}.
\end{itemize}

\section*{Matériel d'accompagnement}

\begin{itemize}
    \item \textbf{Annexes}~: glossaire, outils, papiers de travail types,
          textes, bibliographie, contacts (Annexes~\ref{ann:glossaire}
          à~\ref{ann:contacts}).
    \item \textbf{Livret de correction}~: corrigés des exercices et des
          fiches, réservé au formateur.
    \item \textbf{NDAP SUKULU}~: fiches d'auto-évaluation et laboratoires,
          présentés dans la section qui suit.
\end{itemize}
